% IF-08 and IF-09: actual coefficient kernels and actual m=2 columns.
% Included by contribution.tex. Full author proof reading recorded separately.
\subsection{IF-08: the actual finite theta jet has an exterior-tensor realization}
The finite coefficient map in the source is more general than a single
raising operator. It can nevertheless be realized explicitly and
integrally in exterior tensors, with exactly the same Smith indices.
Here the relevant source declarations are
\texttt{LocalToGlobal/KernelLattices/ThDistC.lean},
\texttt{subD,jetLin,jetMap,jetOp,JetInVV,derivAt0,jetC}
(lines 23--61), and
\texttt{LocalRingCardinality/DataKerI.lean},
\texttt{KM.Gxz,KM.applyAt0,KerEnt.kform} (lines 27--35 and 115--121).
The former is the coefficient map itself; the latter describes its
matrix entries tested against the dual product-minor lattice.

Fix a partition $\lambda$, with column heights $h_1,\ldots,h_q$,
and integers $a,b\geq\ell(\lambda)$. Put
\[
 A_\lambda(r)=\bigotimes_{c=1}^q\Lambda^{h_c}\mathbb Z^r,
 \qquad T_\lambda(r)\subset A_\lambda(r)
 \quad (r=a,b),
\]
where $T_\lambda(r)$ is the divided-root cyclic lattice of IF-04.
Write $\pi_r:A_\lambda(r)\longrightarrow\mathbb Z[X]$ for
multiplication of the corresponding initial-column minors. Use polynomial
matrices with $r$ rows and the common $\ell(\lambda)$ columns. This
makes rectangular substitution meaningful also when $b<a$; the
source \texttt{subD} sends unused source columns $s\geq b$ to zero,
and no such column occurs in these minors. Its cyclic
restriction is the isomorphism
$\Phi_r:T_\lambda(r)\xrightarrow\sim\Delta_\lambda(r)$ of IF-04.
The full tensor image is $\nabla_\lambda(r)$, by the supplied
\texttt{GradedWeylDivisibility/FinalC.lean}, theorem
\texttt{standardMonomialSpan}, whose exact statement is
\texttt{StatementA.lean}, \texttt{StandardMonomialSpan}, lines 221--226.
We use that supplied integral straightening result as background.

Let $Z=(z_{ui})$ be a $b$ by $a$ matrix of independent variables and set
\[
 A_\lambda(Z)=\bigotimes_c\Lambda^{h_c}Z
       =\sum_\beta Z^\beta A_\beta,
 \qquad A_\beta\in\operatorname{Hom}_{\mathbb Z}
       (A_\lambda(a),A_\lambda(b)).
\]
For a characteristic-zero coefficient field $K$ and a prescribed moment
function $J:\mathbb N^{ba}\to K$, define the concrete finite matrix
\[
 B_J=\sum_{|\beta|=|\lambda|}J(\beta)A_\beta.       \tag{IF-08a}
\]
Only finitely many entries of $J$ are used. If
$\Theta(Z)=\sum_\beta\theta_\beta Z^\beta$, the source convention is
\[
 J_\Theta(\beta)=\beta!\,\theta_\beta,
 \qquad \beta!=\prod_{u,i}\beta_{ui}!,              \tag{IF-08b}
\]
including the explicitly prescribed powers of the discriminant-period
generators in \texttt{jetC}. Thus $J$ consists of ordinary derivatives
at zero; it is not an unspecified divided-power normalization.

\paragraph{Statement IF-08.}
The lattices $T_\lambda(r)$ are saturated in $A_\lambda(r)$, and
\[
 \Phi_r:T_\lambda(r)\simeq\Delta_\lambda(r),\qquad
 \Phi_r:T_\lambda(r)^\#\simeq\nabla_\lambda(r)       \tag{IF-08c}
\]
are isometries for the tensor-exterior pairing on the left and the
normalized Fock pairing on the right. The matrix $B_J$ preserves the
rational cyclic spaces. If its moments lie in a subring $R\subset K$,
it carries $T_\lambda(a)\otimes R$ into $T_\lambda(b)\otimes R$.
On those spaces the actual polynomial substitution followed by moment
contraction satisfies
\[
 \Phi_b B_J=\operatorname{jet}_J\Phi_a.             \tag{IF-08d}
\]
This equation is an equality of explicit coefficient matrices. With
source and target lattices as displayed, it preserves every
Smith invariant over $\mathbb Z$ and every elementary-divisor exponent
over a discrete valuation ring. The transpose identity gives the same
statement for the normalized dual lattices.

For a determinant insertion of order $n\geq0$, multiplication by
$\det(X)^n$ is represented on the exterior side by insertion of $n$
copies of the oriented top wedge of $\mathbb Z^a$. Thus the complete
\texttt{jetMap J n}, including its determinant factor, is intertwined
by the same construction. Its rational retraction \texttt{retrVV}
has no effect under the supplied \texttt{JetInVV} hypothesis.

\begin{proof}
We first make the scalar normalization and saturation explicit. Let
\[
 c_\lambda=\langle p_\lambda,p_\lambda\rangle_{\rm Fock}
   =\sum_\alpha ([X^\alpha]p_\lambda)^2\alpha!>0.
                                                               \tag{IF-08e}
\]
This explicit finite sum is an integer independent of the ambient rank $r\geq\ell(\lambda)$.
IF-04 proves $\pi_r^*\pi_r=c_\lambda$ on $T_\lambda(r)_\mathbb Q$,
where the adjoint uses the unnormalized factorial Fock form and the
orthonormal exterior tensor form. The rational image of $\pi_r$ is
already the image of its cyclic subspace, by
\texttt{GradedWeylDivisibility/FinalB.lean}, theorem
\texttt{weylEqNablaQ} (the supplied \texttt{WeylEqNablaQ} statement).
Hence $T_\lambda(r)_\mathbb Q=(\ker\pi_r)^\perp$. On the whole tensor
space we therefore have
\[
 \pi_r^*\pi_r=c_\lambda P_r,\qquad
 P_r=\frac1{c_\lambda}\pi_r^*\pi_r,                 \tag{IF-08f}
\]
where $P_r$ is the orthogonal projection onto the cyclic rational
space. These formulas display, rather than discard, the scalar.

Suppose $t\in T_\lambda(r)_\mathbb Q\cap A_\lambda(r)$.
For every $v\in A_\lambda(r)$, (IF-08f) gives
\[
 \frac{\langle\pi_rt,\pi_rv\rangle_{\rm Fock}}{c_\lambda}
       =\langle t,v\rangle\in\mathbb Z.
\]
The full image $\pi_r(A_\lambda(r))$ is $\nabla_\lambda(r)$.
The supplied normalized duality theorem
\texttt{LocalToGlobal/Columns/CoefficientsModP.lean},
\texttt{weylEqDualNabla} (line 31),
therefore implies $\pi_rt\in\Delta_\lambda(r)$; IF-04 gives
$t\in T_\lambda(r)$. This proves saturation. The other normalized
duality, \texttt{nabla\_eq\_dual\_weyl\_holds} in
\texttt{Columns/CoefficientsModP/WeTransposeWordTransposeWord.lean}
(line 74), and the isometry of IF-04 yield
(IF-08c). Equivalently $P_r(A_\lambda(r))=T_\lambda(r)^\#$;
this also follows by extending an integral linear functional on the
primitive sublattice $T_\lambda(r)$ to the free ambient lattice.
The projection can have denominators, exactly those already present
in this specified dual lattice.

The exterior matrix $A_\lambda(Z)$ has entries
$\prod_c\det Z[I_c,J_c]$. Multiplication of minors gives
\[
 \pi_b A_\lambda(Z)=\rho_Z\pi_a,                  \tag{IF-08g}
\]
where $\rho_Z$ substitutes $X_{is}\mapsto\sum_u Z_{ui}X_{us}$.
For completeness, expand a determinant of these sums by multilinearity
in its rows. Terms with a repeated index vanish by alternation; grouping
all permutations of each increasing index set gives the determinant
$\det Z[I,J]$ times the corresponding minor of $X$. Multiplying these
identities over the columns gives (IF-08g), including all signs.

For a rational specialization of $Z$, (IF-08g) sends $\ker\pi_a$
into $\ker\pi_b$. Its transpose version sends $\ker\pi_b$ into
$\ker\pi_a$, since $A_\lambda(Z)^t=A_\lambda(Z^t)$.
Taking orthogonal complements shows
$A_\lambda(Z)T_\lambda(a)_\mathbb Q\subset
T_\lambda(b)_\mathbb Q$. This holds for every rational matrix $Z$.
Every coefficient $A_\beta$ consequently has the same containment:
apply any linear functional vanishing on the latter space and use
that a polynomial vanishing at all rational tuples is zero.
The coefficients $A_\beta$ are integer matrices; the saturation just
proved implies $A_\beta T_\lambda(a)\subset T_\lambda(b)$.
This proves the claimed integrality for arbitrary $R$ and $J$.

Apply the moment functional $Z^\beta\mapsto J(\beta)$ to (IF-08g).
It gives (IF-08d). For $J=J_\Theta$, direct differentiation of a
monomial proves (IF-08b); this is exactly the factorial in
\texttt{KM.applyAt0}. The test-minor version is
\[
 \operatorname{applyAt0}(\Theta,G_I(XZ))
   =\sum_J \Bigl(\sum_\beta\beta!\theta_\beta
                [Z^\beta]\prod_c\det Z[J_c,I_c]\Bigr)G_J(X),
                                                               \tag{IF-08h}
\]
with the direction of $Z$ transposed to match the source's test
convention. This is the same matrix, paired with the dual lattice.
Under (IF-08c), normalized Fock testing is ordinary tensor-exterior
testing. The $c_\lambda$ in (IF-08f) has canceled because the Fock
form is divided by that same $c_\lambda$; replacing normalized Fock
by unnormalized Fock would retain that factor.

Finally insert $n$ full-height columns in $A_\lambda(a)$. Their top
wedges have norm one and transform by $\det^n$, while every row root
annihilates them. Their cyclic sublattice is therefore exactly the
cyclic lattice of the shape $\lambda+n(1^a)$, and their minor product
is $\det(X)^n$. Thus determinant insertion also has no unspecified
scalar. In integral bases, (IF-08d) and its determinant variant are
conjugations by unimodular source and target matrices. Their images,
cokernels, saturated images and local lengths are correspondingly
isomorphic. This proves the assertions about Smith invariants,
including for rectangular maps and rank-deficient images, where the
finite index means the index in the saturation of the image.
\end{proof}

\paragraph{Assembly of the actual finite kernel.}
This is not merely a highest-weight-root comparison. For every
supplied theta jet $J_{j,c',c}$, (IF-08a) specifies its exterior-tensor
matrix. The complete block matrix retains the source's actual scalar
factors:
\[
 \widehat K_j(c',c)=\frac{\rho_j(c)\rho'_j(c')}{|\Gamma_{j,c}|}
        B_{J_{j,c',c}}\,D_{n_j}.                  \tag{IF-08i}
\]
Here $D_{n_j}$ inserts the determinant factors, and the target
one-dimensional determinant character is retained as a separate
rank-one lattice. The scalar entries and the jets include precisely
the period generators and frame changes specified in the source.
Formula (IF-08i) is the tensor version of
\texttt{SharpRestriction/GlobalSchurBounds/GL3Theta.lean},
\texttt{liftTie,lift\_formula} (lines 575--596), and
\texttt{LocalRingCardinality/DataKernels.lean},
\texttt{wt,sharp,divDe,Kmap,K} (lines 52--68).

The coefficient isomorphisms commute with the relevant group actions;
hence they transport invariant submodules, local lattice tests and
finite direct sums over lattice classes. For an actual supplied
finite global matrix they give an isomorphic kernel matrix, with all
its stabilizer denominators still displayed. The test condition
\texttt{KerEnt.lat\_S} is pairing integrality against
\texttt{KM.nablaX}; its exterior version is integrality against
$T_\lambda^\#$, not against a second copy of $T_\lambda$.
When the source's local integrality assumptions hold, the transported
matrix has exactly the same integral elementary divisors. No new
integrality assumption is silently supplied here.

The precise global integral lattice is also identifiable in the source.
\texttt{Columns/Record.lean}, \texttt{stepH} (lines 1286--1287),
uses \texttt{intDel}; \texttt{Columns/CoefCol.lean},
\texttt{roundPk} (lines 59--64), uses \texttt{Del=delAtr} and
\texttt{Nab=nabAtr}. Thus the integral global column forms use the
Weyl lattice; the dual Weyl lattice supplies its testing lattice.
In \texttt{Divisibility/CoefMod.lean},
\texttt{transLat,delP,delK} (lines 196--211), the lattice is explicitly
the intersection of the inverse images of standard local lattices
under \texttt{actP} applied to $B^{-1}$ for local frames $B$.
Since \texttt{actP} is determinant twist times the transposed
polynomial row action (lines 69--76), its exterior action is
\[
 R_\lambda(B)=\prod_\tau\det(B_\tau)^{s_\tau}
       \bigotimes_{\tau,c}\Lambda^{h_{\tau,c}}B_\tau.
                                                               \tag{IF-08j}
\]
The round realization uses the source's inverse-transpose convention
\texttt{rhoVr}; its \texttt{delAtr} is formed from the dual frame
lattice (\texttt{CoefMod.lean}, lines 320--330). Applying exactly that
same convention to (IF-08j) transports the global intersection and
all its localizations; no free global basis of a nonprincipal
coefficient ideal is presumed.

In the following formula $B_{j,c}$ means the frame of the coefficient
lattice used by \texttt{transLat}. For an actual round column whose
geometric frame is $C_{j,c}$, this means $B_{j,c}=C_{j,c}^{-t}$,
with the dual-frame convention just specified; one must not insert
$C_{j,c}$ directly into $R_j(B_{j,c})$. For chosen local bases at a
prime $w$, the transported block matrix is therefore the explicit matrix
\[
 R_{j+1}(B_{j+1,c'})^{-1}
 \frac{\rho_j(c)\rho'_j(c')}{|\Gamma_{j,c}|}
 \Bigl(\sum_{\beta}J_{j,c',c}(\beta)A_{\beta}\Bigr)
 D_{n_j} R_j(B_{j,c}).                              \tag{IF-08k}
\]
One sums these blocks over the actual finite class representatives,
then restricts to the invariant and level submodules defining forms.
This is the double-coset sum \texttt{kerOpM}; its scalar analogue is
\texttt{thetaLiftAMF} in \texttt{ThetaLiftIndependence/Statement.lean},
lines 80--84. In a local field the scalar valuation in (IF-08k) is
\[
 v_w(\rho_j(c))+v_w(\rho'_j(c'))-v_w(|\Gamma_{j,c}|),
\]
and a jet additionally contributes the recorded period-generator
valuation $\sum_\tau |\beta_\tau|v_w(c_{D,\tau})$ to its ordinary
Taylor derivative valuation. They have not been dropped or replaced
by root factorials. For a resulting local matrix $M$ of rank $r$,
put $s_0=0$ and
\[
 s_k=\min_{|I|=|J|=k}v_w(\det M_{I,J}),\qquad 1\leq k\leq r.
\]
Its elementary-divisor exponents are $s_k-s_{k-1}$; its image index
in the saturated image, when $M$ is integral, has length $s_r$.
These statements follow directly from the Smith normal form over a
DVR, and the formulas are unchanged by (IF-08c). This gives an exact
finite algorithm for the actual global transport, including every
local denominator. A terminal period functional is transported by
the same dual isomorphism, so its value on the transported chain is
identical, not merely bounded by the original value.

The result does not identify (IF-08i) with a single $L$, $L^{(r)}$,
or a polynomial in one Lefschetz pair. The actual moments,
class numbers, characters and stabilizers remain part of the matrix.
It constructs the requested finite-kernel intertwiner and proves
index transport; equality with a separately defined Lefschetz order
index would still require an additional operator identity. IF-03 and
IF-05 show why discarding the operator and lattice distinctions gives
false numerical equalities.

\subsection{IF-09: the actual two-step column chain and a fixed kernel test}
The polynomial coefficient weights of the definite column chain must
be distinguished from the auxiliary $a^+$ row weights in IF-05.
For the actual \texttt{numEx} datum and its all-false word, define
\[
 u=(113/2,47/2,39/2,31/2,23/2,15/2,7/2,-1/2),\qquad
 \lambda_0=(53,21,18,15,12,9,6,3).
\]
The unshifted weights from \texttt{MemberForms.lam\_eq} are
\[
 \lambda_1^N=(\lambda_0,-88^{[16]}),\qquad
 \lambda_2^N=(\lambda_0,-88^{[16]},-90^{[16]}).
\]
The actual \texttt{colVal} uses \texttt{Wts.lamTrue}, which subtracts
$8j$ from every entry at stage $j$. Thus the determinant exponents of
its three polynomial realizations are
\[
 (s_0,s_1,s_2)=(3,-96,-106),
\]
and their polynomial partitions are
\[
\begin{split}
 \alpha_0&=(50,18,15,12,9,6,3,0),\\
 \alpha_1&=(141,109,106,103,100,97,94,91,0^{[16]}),\\
 \alpha_2&=(143,111,108,105,102,99,96,93,2^{[16]},0^{[16]}).
\end{split}                                                     \tag{IF-09a}
\]
Here $x^{[r]}$ means $r$ repeated entries. These formulas follow by
substitution in \texttt{DataCoeff.hcToWeightQ},
\texttt{Params.wG}, \texttt{MemberForms.lam\_eq}, and
\texttt{Columns/AbelianHostFieldConditions.lean},
\texttt{Wts.lamN, Wts.lamTrue} (lines 462--465).
\texttt{Columns/CoefCol.lean}, \texttt{colVal} (lines 89--91), and
\texttt{Divisibility/ColPins.lean}, \texttt{lamOfCol} (lines 21--26),
pin this distinction to the actual column map.

The kernel diagrams, before appending target zeros, are
\[
 \nu_1=(141,109,106,103,100,97,94,91),\qquad
 \nu_2=(143,111,108,105,102,99,96,93,2^{[16]}).       \tag{IF-09b}
\]
They are exactly the output of \texttt{NumDatum.nu} and
\texttt{DV.nuRow}; they are also
$\alpha_0+91(1^8)$ and $\alpha_1+2(1^{24})$.
The source has
$\mathrm{sSharp}=(88,98)$ and $\mathrm{sDiv}=(96,106)$.
The identities
$s_i+\mathrm{sSharp}_i=n_i$ and
$s_{i+1}=-\mathrm{sDiv}_i$ give
\[
 (n_0,n_1)=(91,2).                                  \tag{IF-09c}
\]
They are precisely \texttt{KerEntTied.WtTie} in
\texttt{KernelLattices/JetFieldLeKR.lean}, lines 108--111.
Consequently the actual finite coefficient chain of IF-08 is
\[
 T_{\alpha_0}(8)\xrightarrow{D_{91}}
 T_{\nu_1}(8)\xrightarrow{B_{J_1}}
 T_{\alpha_1}(24)\xrightarrow{D_2}
 T_{\nu_2}(24)\xrightarrow{B_{J_2}}
 T_{\alpha_2}(40),                                  \tag{IF-09d}
\]
The two determinant degree increments are $8\cdot91=728$ and
$24\cdot2=48$. Thus two lifts do not represent $L^2$.
The chain retains the determinant characters $s_i$ and the scalar block
factors (IF-08i). This specifies the whole finite polynomial kernel,
not just a simple-root summand.

The actual column simple-root gaps include $32,3,91,2$ (zero gaps are
trivial). Their top-to-second-descendant invariants are
\[
\begin{array}{c|c|c|c}
 d&F^2\text{ on Weyl}&F^2\text{ on normalized dual}
       &F^{(2)}\text{ on normalized dual}\\\hline
 2&2&2&1\\
 3&2&6&3\\
 32&2&992&496=2^4\cdot31\\
 91&2&8190&4095=3^2\cdot5\cdot7\cdot13
\end{array}.                                                   \tag{IF-09e}
\]
On the Weyl lattice every displayed $F^{(2)}$ invariant is one, by
IF-04. Thus the dual-lattice obstruction is present in the actual
column coefficients themselves. The auxiliary row gap $90$ in IF-05
is a separate, correctly labeled test.

The source's \texttt{numEx} is a numerical datum, not a full
\texttt{NumFacts} witness. The repaired \texttt{numBig} is given in
\texttt{NumbersLink.lean}, \texttt{numBig} (line 803) and
\texttt{numFactsBig} (line 1260). It retains the row parameters and
changes $u_7$ from $-1/2$ to $-7/2$ at one specified real place, besides
changing the size parameters. At that place $s_0=0$, $n_0=88$;
the last entries $91,93$ in (IF-09b) become $88,90$ respectively.
The other entries and the second determinant insertion $n_1=2$
are unchanged. The additional column gaps are $6,88$, with
$\binom62=15$ and $\binom{88}2=3828=2^2\cdot3\cdot11\cdot29$.
The normalized dual obstruction therefore persists for this valid
numerical-facts instance. It does not depend on treating
\texttt{numEx} as a nonexistent full global witness.

There is also a precise host obstruction to regarding
\texttt{numBig} itself as a full global \texttt{D7} record.
\texttt{NumbersLink.numBig} retains \texttt{numEx.H};
\texttt{NumbersLemmas.hostEx\_q,hostEx\_d} give $q=11,d=7$.
But \texttt{DataHost.lean}, \texttt{Host.qd\_big} (line 52), requires
$4tr+1<q,d$. Here $t=15,r=9$, so $4tr+1=541$, contradicting both
inequalities. Thus its \texttt{NumFacts} theorem certifies numerical
conditions only; it cannot instantiate that full global host with
these primes. The polynomial representations and kernel tests above
are nevertheless fully specified and valid on their stated coefficient
lattices.

\paragraph{An exact nonzero kernel test.}
For the rectangular coordinate inclusion $A:\mathbb Z^a\hookrightarrow
\mathbb Z^b$, define
\[
 \Theta_A(Z)=\exp\bigl(\sum_{u,i}A_{ui}z_{ui}\bigr),\qquad
 [Z^\beta]\Theta_A=A^\beta/\beta!.
\]
This is the precise series \texttt{ThE} in
\texttt{SpanPowerMembership/NZRowV.lean}, lines 77--79.
Its factorial cancellation is the supplied theorem
\texttt{fac\_ThE}, lines 89--97; its kernel evaluation is
\texttt{eval\_applyAt0,eval\_Gxz}, lines 99--146.
For this specified series $J(\beta)=A^\beta$, so
$B_J=A_\lambda(A)$ by (IF-08a). Each map in (IF-09d) is then an
inclusion with saturated image. Indeed it carries its ambient tensor
basis injectively into the larger ambient tensor basis; saturation of
the cyclic lattices was proved in IF-08. The determinant insertion
has the same property, since it tensors with a primitive norm-one
vector. Their composite is primitive for the same ambient basis
reason. Its nonzero Smith factors are all one.

Taking the highest vector at the beginning and the normalized dual
highest functional at the end gives period exactly $1$ after both
arrows. This assertion is fully fixed-data: the kernels, lattices,
vectors, determinant insertions and functional have been specified.
On the highest diagonal monomial the derivative factor is
$\prod_i\nu_{j,i}!$ and the coefficient of $\Theta_A$ is its inverse,
so the unit value is an explicit factorial cancellation, not an
assumption that the two normalizations agree.

\paragraph{A compatible common Lefschetz action on this fixed family.}
There is nevertheless a positive integral intertwining result for the
same specified kernel chain. Fix one simple root inside the first eight
rows, for example rows $2,3$. Every determinant insertion and coordinate
inclusion in (IF-09d) commutes with both of its root operators $E_i,F_i$:
$E_i$ and $F_i$ kill a full determinant vector, and the coordinate
inclusion is the identity on those two rows. These identities hold on
the entire cyclic coefficient lattice, not just the highest root string.

For each coefficient lattice apply the particle-hole embedding of IF-04
into $V_{N_j}$, with $N_j=(r_j-1)q_j$ and $q_j$ its number of columns.
Let $N=\max_jN_j$, and pad the embedding with a singleton $e$ in each
of the $N-N_j$ extra symplectic pairs. Singleton padding is an integral
isometry with primitive image and is annihilated by both the raising
and lowering on each extra pair. It therefore intertwines the same
$L,D$ on $V_N$. Primitivity of $T_\lambda$ in its column ambient lattice
(IF-08), followed by the coordinate-basis particle-hole embedding and
padding, makes each image a primitive $L,D$-stable sublattice of $V_N$.
Thus the fixed \texttt{ThE} arrows are specified integral isometric maps
between those sublattices, commuting with one common Lefschetz action,
and their nonzero Smith factors remain one. For (IF-09d), one may take
$N=39\cdot143=5577$ at each embedding factor. Tensoring the eight
factors gives the diagonal root action on $V_{8N}=V_{44616}$, with
the same integral and primitivity statements.

The exterior degree of a vector of simple-root weight $h$ is
$N_j-h$ before padding and $N-h$ after padding. Consequently these
intertwining maps have degree zero. They are not applications of $L$;
that stronger assertion is excluded below. No integral equivariant
extension of the arrows to all of $V_N$ is asserted. Nor is algebraicity
of these particular particle-hole image sublattices asserted; IF-10
provides a separate geometric coefficient realization with its stated
character normalization.

For an arbitrary moment block $B$ between the same standard coefficient
lattices, the exact additional common-root conditions are the finite
matrix equations
\[
 E_i' B=B E_i,\qquad F_i' B=B F_i.                 \tag{IF-09f}
\]
They suffice for all divided powers, by equality in characteristic zero
and their integral preservation. For the actual global form spaces one
must also check that the proposed root action preserves their invariant
and level submodules and their local lattices. These are explicit
operator/lattice conditions after (IF-08k), and are not consequences of
mere theta-kernel rationality. No such common-root invariance is assumed
for the unspecified global moments.

\paragraph{A fixed-normalization obstruction to two ordinary Lefschetz arrows.}
This particular two-step kernel chain cannot be intertwined with two
successive applications of the same fixed ordinary $L$ on one full
integral symplectic exterior lattice by integral embeddings whose
final target image is saturated. Indeed its highest vector is
sent to a primitive highest vector. A saturated integral embedding
preserves primitivity. But $L^2=2L^{(2)}$ on every such exterior
lattice, so the image of $L^2$ is contained in $2V_n$ and contains no
primitive vector. This is a contradiction. Even without saturation,
an isometric integral embedding for the fixed normalized pairings is
impossible: the resulting highest vector has squared norm one, whereas
a nonzero vector in $2V_n$ has squared norm at least four. The full
exterior monomial basis is orthonormal, so both assertions follow
immediately by inspecting its integer coefficients. This proves a
fixed-data failure of the ordinary-$L^2$ normalization, rather than a
failure caused by leaving the test vector unspecified. The same argument
also excludes two positive divided powers of that same fixed $L$:
for $a,b>0$,
\[
 L^{(a)}L^{(b)}=\binom{a+b}{a}L^{(a+b)},
\]
and the integer $\binom{a+b}{a}>1$ divides every coordinate of the
composite image. Thus a primitive or norm-one composite cannot occur.
This leaves tensor-functor realizations such as IF-09d, or operators
belonging to different specified Lefschetz pairs, outside the
obstruction.

These kernels are a concrete algebraic test family for the exact
source operator formula; the source also builds \texttt{KerEnt}
instances with \texttt{ThE} in \texttt{NZRowW.lean},
\texttt{levelRecordTwistEntries}, and \texttt{NZRowZ.lean},
\texttt{fin\_ent}. We do not assert that this exponential test family
is the globally automorphic theta function in Alp\"oge's construction,
or that it satisfies all the additional global theta hypotheses.
For that actual global function, (IF-08i) is still an explicit
intertwining formula in its actual finite jets; computing its Smith
numbers requires those finite jet values and class data. An
existential theta construction or a lower bound does not determine
those numbers. The finite-kernel realization has therefore been
constructed, while the distinct claim that every actual global index
equals one universal Lefschetz order index has not been proved.

\paragraph{Schemes, $\lambda$-operations, and realizations.}
Every $A_\lambda(Z)$ above is a polynomial matrix over
$\operatorname{Spec}\mathbb Z$, assembled from the exterior-power
operations $\Lambda^{h_c}$ and tensor products. Its class in $K_0$
is the corresponding product of $\lambda^{h_c}$ classes. The
saturated cyclic lattice and its explicit dual retain information
which $K_0$ alone forgets: a finite-index inclusion has the same
rational character and can have different integral cokernel.
Base change to $\mathbb Z_\ell$ preserves the displayed integral
intertwiners and carries a Smith factor $d$ to local length
$v_\ell(d)$. Tensoring with $\mathbb Q_\ell$ forgets that length.
When these spaces are realized in the cohomology of an abelian
variety, their tensor and exterior operations are motivic tensor
operations. The formal coefficient identification alone supplies no
new algebraic correspondence representing the theta moment matrix
or the Hodge sign operator. Such an assertion would require the
separate algebraicity comparison under discussion in Q6g and the
main question.
